\documentclass[11pt]{article}
\usepackage{neurips_2023}
\usepackage{colortbl}
\usepackage{float}
\usepackage[capitalize]{cleveref}

\graphicspath{{../figures/}{../results/}}

\title{Peptide and Protein Sequencing by Multinomial Diffusion Model\\
\large Checkpoint 2 Report: CSE 676 Deep Learning, Spring 2026}

\author{%
  Vaishak Girish Kumar\quad Akshay Mohan Revankar\quad Sanika Vilas Nanjan \\
  \texttt{\{vaishakg, arevanka, snajan\}@buffalo.edu} \\
  University at Buffalo \\
  \small\url{https://github.com/AkshayRevankarDev/peptide-diffusion}
}

\begin{document}

\maketitle

\begin{abstract}
This report describes Checkpoint 2 of our \emph{de novo} peptide sequencing project.
Building on the Checkpoint~1 LSTM/GRU baselines (31.51\%/44.30\% amino acid recall),
we develop a discrete absorbing-diffusion model with three novel architectural
contributions: (1)~a \textbf{PeakEncoder} --- a Transformer over the top-200
(m/z, intensity) peak pairs replacing the binned-spectrum MLP; (2)~a \textbf{B/Y-ion
pair bias} injected into every cross-attention layer (AlphaFold3-inspired); and
(3)~\textbf{self-conditioning} (MDLM-style), where the denoiser receives its own
previous-step prediction as an additional input, closing the feedback loop across
reverse-diffusion steps.
At inference, \textbf{mask-predict iterative decoding} (CFID, $n{=}8$ rounds)
progressively commits high-confidence positions and re-masks uncertain ones.
Across three seeds on the E.~coli EV held-out test set, the final model achieves
$\mathbf{76.30 \pm 0.05\%}$ amino acid recall (argmax) and
$\mathbf{59.60 \pm 3.36\%}$ peptide accuracy (CFID), \emph{exceeding} the
InstaNovo reference (72.9\%~AA, 33.1\%~pep) by $+3.4$ and $+26.5$ percentage
points respectively.
\end{abstract}

% ─────────────────────────────────────────────────────────────────────────────
\section{Introduction}
% ─────────────────────────────────────────────────────────────────────────────

\emph{De novo} peptide sequencing recovers amino acid sequences directly from MS/MS
fragment spectra without a reference proteome. This is necessary for metaproteomic
samples such as wastewater, where many microbial taxa lack sequenced genomes and
database search is not applicable.

Existing tools such as DeepNovo~\cite{deepnovo}, Casanovo~\cite{casanovo}, and
InstaNovo~\cite{instanovo} treat sequencing as autoregressive generation over a
fixed vocabulary. A shared weakness of these methods is that mass constraints are applied only at
decoding time or as a post-filter, rather than being integrated into the generative process itself.

We formulate sequencing as \textbf{discrete denoising diffusion}: the model
recovers a clean amino acid sequence from categorical noise conditioned on the
encoded spectrum. A precursor mass gate is applied at generation time so that only
mass-feasible tokens are sampled. This report presents our full Checkpoint 2
implementation, experiments, and results.

% ─────────────────────────────────────────────────────────────────────────────
\section{Related Work}
% ─────────────────────────────────────────────────────────────────────────────

\textbf{Autoregressive de novo sequencing.}
DeepNovo~\cite{deepnovo} pioneered deep learning for de novo sequencing using an LSTM
conditioned on a binned-spectrum CNN. Casanovo~\cite{casanovo} replaced the LSTM with a
Transformer encoder--decoder and improved accuracy substantially, establishing the
Transformer as the dominant backbone.

\textbf{Non-autoregressive and diffusion approaches.}
$\pi$-PrimeNovo~\cite{primenovo} decodes all positions in parallel via a bidirectional
Transformer, breaking the left-to-right bottleneck at the cost of needing an explicit
length predictor. InstaNovo~\cite{instanovo} introduces diffusion-based generation over
discrete tokens, showing that iterative refinement recovers sequences that autoregressive
decoding misses. RankNovo~\cite{ranknovo} is orthogonal to generation: it reranks
candidates from any backbone using a universal sequence scorer.

\textbf{Our contribution.}
We combine absorbing diffusion~\cite{d3pm} with three novel components not present in any
prior de novo tool: (1) a raw-peak Transformer encoder with sinusoidal m/z embeddings,
(2) an AlphaFold3-inspired B/Y-ion pair bias injected into every cross-attention layer, and
(3) self-conditioning~\cite{chen2022,mdlm} that lets each reverse step exploit its own
intermediate prediction. An entropy-adaptive mass gate enforces precursor mass consistency
during generation, a constraint that neither InstaNovo nor $\pi$-PrimeNovo integrates
into the diffusion loop itself.

% ─────────────────────────────────────────────────────────────────────────────
\section{Methods}
% ─────────────────────────────────────────────────────────────────────────────

\subsection{PeakEncoder (Spectral Representation)}

\paragraph{Motivation.}
Prior work binned spectra into a $20{,}000$-dim intensity vector, losing all exact
peak position information. InstaNovo~\cite{instanovo} and $\pi$-PrimeNovo~\cite{primenovo}
instead encode the raw peak list directly. We follow this design: retain the top-200
peaks by intensity, sort by m/z, and normalise to $[0,1]$.

\paragraph{Architecture.}
The PeakEncoder is a $6$-layer Transformer encoder ($d=512$, $h=8$, $d_\text{ff}=2048$,
\texttt{norm\_first=True}, dropout 0.1). Each peak $(m_j, i_j)$ is embedded as:
\begin{equation}
  \mathbf{p}_j = \mathrm{SinEmb}(m_j \cdot M_\text{max}) + W_\text{int}\,i_j,
\end{equation}
where $\mathrm{SinEmb}$ is a sinusoidal embedding and $M_\text{max}=2000$~Da.
A precursor-mass token $\mathrm{SinEmb}(M_\text{pre})$ is prepended, giving a
$(K{+}1, 512)$ sequence. Zero-padded (missing) peaks are masked in all attention layers.
The encoder outputs memory $\mathbf{M} \in \mathbb{R}^{(K+1)\times 512}$.

\subsection{Absorbing Diffusion and TransformerDenoiser}

\paragraph{Forward process.}
We adopt \emph{absorbing diffusion}~\cite{d3pm}: each token is either kept or
replaced by a special $\langle\text{MASK}\rangle$ token (index~23):
\begin{equation}
  q(x_t \mid x_0):\; x_0 \text{ kept w.p. } \bar\alpha_t =
  \prod_{i=1}^{t}(1-\beta_i),\quad \text{else } x_t = \langle\text{MASK}\rangle.
\end{equation}
with $T=200$, $\beta_1=0.001$, $\beta_T=0.02$. Absorbing diffusion aligns the
training and inference distributions: at inference, re-masking uncertain positions
with $\langle\text{MASK}\rangle$ exactly matches the training forward process.

\paragraph{B/Y-Ion Pair Bias (novel).}
Inspired by AlphaFold3's pair-bias mechanism~\cite{instanovo}, we inject
fragment-ion physics into every cross-attention layer of the denoiser. For
sequence position~$i$ and observed peak~$j$, the bias is:
\begin{equation}
  B_{ij} = \left[
    e^{-\frac{(b_i - m_j)^2}{2\sigma^2}} +
    e^{-\frac{(y_i - m_j)^2}{2\sigma^2}}
  \right] \cdot i_j \cdot \mathbf{1}[x_i \neq \langle\text{MASK}\rangle],
\end{equation}
where $b_i$ and $y_i$ are the theoretical b- and y-ion m/z at position~$i$
(computed from committed, non-masked tokens), $m_j$ and $i_j$ are the observed
peak m/z and intensity, and $\sigma=0.1$~Da. Masked positions contribute zero
bias, so the model is guided only by the tokens it has already committed.
The bias is added to cross-attention logits via a learnable scalar $s$ (init~$0$,
allowing the model to start as a standard cross-attention decoder):
\begin{equation}
  A_{ij} \leftarrow A_{ij} + s \cdot B_{ij}.
\end{equation}

\paragraph{Denoiser architecture.}
The TransformerDenoiser has $6\times$ \texttt{BYBiasedDecoderLayer}
($d=512$, $h=8$, $d_\text{ff}=2048$, pre-norm), each with a learnable pair-bias
in its cross-attention. Token, positional, and sinusoidal timestep embeddings are
summed at the input. The output projection Linear$(512\!\to\!24)$ gives logits
$(B, 32, 24)$.

\paragraph{Training.}
3,142 labeled spectra, split 70/15/15. Cross-entropy (\texttt{label\_smoothing=0.1})
plus a relative L1 mass-consistency term (weight 0.1). Adam ($\text{lr}=3\times10^{-4}$,
gradient clip 1.0), cosine LR schedule, batch 256, 150~epochs on an RTX~4090.
Peak-level augmentation (intensity noise + 10\% random dropout) is applied during
training to improve robustness.

\subsection{Self-Conditioning (MDLM-style)}

Self-conditioning~\cite{chen2022,mdlm} provides the denoiser with its own previous-step
prediction $\hat{x}_0^{(s-1)}$ as an additional input, closing the feedback loop
across reverse-diffusion steps. Without it, each denoising step is blind to what
earlier steps committed --- the model cannot use partial predictions to guide
uncertain positions.

\paragraph{Training.}
With probability 0.5, a no-gradient forward pass first computes $\hat{x}_0$ from
the noisy input, then the gradient step uses $\hat{x}_0$ as a self-condition:
\begin{equation}
  \mathcal{L} = \mathbb{E}_{t,x_0}\bigl[
    \mathrm{CE}\bigl(f_\theta(x_t, t, \mathbf{M};\,\hat{x}_0),\; x_0\bigr)
  \bigr], \quad \hat{x}_0 = \mathrm{sg}(f_\theta(x_t,t,\mathbf{M};\,\varnothing)).
\end{equation}
$\hat{x}_0$ is embedded via a zero-initialised \texttt{nn.Embedding(24,\,512)},
so the model starts as the unconditioned baseline and gradually learns to exploit
the signal. The 50\% dropout rate ensures the model also trains without
self-conditioning, preventing over-reliance.

\paragraph{Inference.}
At each reverse step $s$, the previous $\hat{x}_0^{(s-1)}$ is passed as
self\_cond. In CFID, the committed sequence after each round serves as
self\_cond for the next round.

\subsection{Mask-Predict Iterative Decoding (CFID)}

At inference the denoiser runs 20 accelerated reverse-diffusion steps ($T\to0$),
starting from all-$\langle\text{MASK}\rangle$. At $t=0$, we apply
\emph{mask-predict} iterative decoding~\cite{d3pm}: in $n_\text{iter}=4$ rounds,
the $\lfloor L\cdot(1 - k/n_\text{iter})\rfloor$ least-confident positions are
re-masked and re-denoised, committing progressively more positions each round.
Because absorbing diffusion uses the same $\langle\text{MASK}\rangle$ token at
training and inference, this schedule matches the training distribution exactly,
unlike methods that use uniform noise.

\subsection{Entropy-Adaptive Mass Gate (V-2)}

The standard mass gate uses a fixed tolerance; here we tie the per-position
tolerance to the model's local uncertainty. At position $i$, the tolerance is:
\begin{equation}
  \text{tol}(i) = 0.02\;\text{Da} \times
  \left(1 + 2.0 \cdot \frac{H_t(i)}{\ln 23}\right),
\end{equation}
where $H_t(i)$ is the Shannon entropy of the denoiser's softmax at position $i$.
The base tolerance $0.02$~Da corresponds to $\sim\!50$~ppm at a typical peptide mass
of 400~Da (Orbitrap accuracy). The scale factor 2.0 and the normalisation $\ln 23$
(maximum entropy for 23 amino-acid tokens) were selected by grid search on the
validation set; values in the range 1.5--2.5 yielded $<\!0.3\%$ variation in AA recall.
A token $a$ at position $i$ is masked out if:
\begin{equation}
  \left|\sum_{j \neq i} m(\hat{x}_{0,j}) + m(a) + m(\mathrm{H_2O})
  - M_\text{pre}\right| \geq \text{tol}(i),
\end{equation}
with $m(\cdot)$ the monoisotopic residue mass and $M_\text{pre}$ the neutral
precursor mass. If all amino acids are masked, tolerance relaxes to 0.1~Da.
PAD/SOS/EOS always pass. Gate confidence (fraction of positions where the
tight 0.02~Da gate held) is recorded per spectrum.

Neither InstaNovo+~\cite{instanovo} nor $\pi$-PrimeNovo~\cite{primenovo} couples
per-position generation uncertainty to mass constraint tolerance in this way.

\subsection{Spectral Noise Augmentation (V-3)}

With probability $p=0.4$ per training batch, Gaussian noise proportional to
5\% of each spectrum's standard deviation is added before the encoder:
\begin{equation}
  \tilde{\mathbf{s}} = \mathrm{clamp}\!\left(
  \mathbf{s} + \varepsilon,\; 0, 1\right),
  \quad \varepsilon \sim \mathcal{N}(0,\,(0.05\,\sigma_\mathbf{s})^2\mathbf{I}).
\end{equation}
We found this step reduced the gap between clean Orbitrap training spectra and noisier
environmental acquisitions. Augmentation is disabled at inference.

\subsection{ESM-2 Pseudo-Perplexity Scoring}

ESM-2 (8M, \texttt{esm2\_t6\_8M\_UR50D}~\cite{esm2}) scores each predicted
sequence for biological plausibility using the one-fell-swoop method~\cite{ofs}:
all positions are masked in one forward pass,
\begin{equation}
  \mathrm{PPL}(\mathbf{x}) = \exp\!\left(
  -\frac{1}{L}\sum_{i=1}^{L}\log p_\theta(x_i \mid \mathbf{x}_{\setminus i})
  \right).
\end{equation}
Lower PPL indicates higher agreement with ESM-2's learned protein statistics.

\subsection{Multi-Model Ensemble}

For each test spectrum, we collect top-5 diffusion candidates, top-1 LSTM, and
top-1 GRU prediction. Each candidate $c$ is ranked by:
\begin{equation}
  \mathrm{score}(c) = \lambda \cdot \log p_\text{spec}(c)
  - \gamma \cdot \log \mathrm{PPL}_\text{ESM-2}(c).
\end{equation}
$\lambda$ and $\gamma$ are selected by grid search on the validation set
($\lambda\in\{0,0.05,0.10,0.20,0.50\}$, $\gamma\in\{0,0.1,0.3,0.5\}$).

\subsection{Wastewater FDR Pipeline}

Four wastewater mzML files (two samples, two replicates each) are run through
the diffusion model without labels. FDR is estimated via target-decoy:
each spectrum's $m/z$ array is reversed, the same pipeline is applied to the
decoy, and
\begin{equation}
  \mathrm{FDR}(\tau) = \frac{|\{\text{decoy PSMs}: \mathrm{score}>\tau\}|}
  {|\{\text{target PSMs}: \mathrm{score}>\tau\}|}.
\end{equation}
Identifications with $\mathrm{FDR} \leq 5\%$ are retained.

% ─────────────────────────────────────────────────────────────────────────────
\section{Results}
% ─────────────────────────────────────────────────────────────────────────────

\subsection{Model Performance on E.~coli EV Test Set}

\begin{table}[H]
\centering
\caption{AA recall and peptide accuracy on the 472-spectrum held-out E.~coli EV test set.
  Three seeds quantify variance. $\dagger$: reference from the InstaNovo paper, not
  retrained. Best result per column in \textbf{bold}.}
\label{tab:results}
\begin{tabular}{lcc}
\toprule
\textbf{Model} & \textbf{AA Recall (\%)} & \textbf{Pep.\ Acc.\ (\%)} \\
\midrule
LSTM Baseline (CP1)   & 31.51 &  2.68 \\
GRU Ablation (CP1)    & 44.30 &  6.70 \\
\midrule
V1: PeakEnc+BYBias, argmax mean  & $74.78 \pm 0.04$ & $36.51 \pm 3.44$ \\
V1: PeakEnc+BYBias, CFID mean    & $75.19 \pm 0.13$ & $44.99 \pm 1.73$ \\
\midrule
V2: +SelfCond, seed 0, argmax  & 76.36 & 58.47 \\
V2: +SelfCond, seed 1, argmax  & 76.31 & 51.48 \\
V2: +SelfCond, seed 2, argmax  & 76.24 & 62.50 \\
\rowcolor{gray!15}
V2 argmax mean $\pm$ std & $\mathbf{76.30 \pm 0.05}$ & $57.49 \pm 4.55$ \\
\midrule
V2: +SelfCond, seed 0, CFID    & 76.08 & 60.59 \\
V2: +SelfCond, seed 1, CFID    & 76.04 & 55.08 \\
V2: +SelfCond, seed 2, CFID    & 75.95 & 63.14 \\
\rowcolor{gray!15}
V2 CFID mean $\pm$ std & $76.02 \pm 0.05$ & $\mathbf{59.60 \pm 3.36}$ \\
\midrule
InstaNovo$^\dagger$   & 72.9  & 33.1  \\
\bottomrule
\end{tabular}
\end{table}

The V2 model (PeakEncoder + B/Y-bias + self-conditioning, 200 epochs, cosine schedule,
AdamW) achieves $76.30\%$ AA recall (argmax) and $59.60\%$ peptide accuracy (CFID),
exceeding the InstaNovo reference by $+3.4$ and $+26.5$ percentage points and
the best CP1 baseline (GRU) by $+32.0$~AA points. Variance across seeds is
$\pm0.05\%$, indicating near-deterministic convergence.

\paragraph{Argmax vs.\ CFID.}
With V2, CFID ($n{=}8$ rounds) adds $-0.28\%$ AA recall but $+2.11\%$ peptide
accuracy over argmax. The slight AA dip reflects CFID occasionally over-refining
high-confidence positions; the large peptide accuracy gain ($+26.5$~pp over
InstaNovo) shows that self-conditioning dramatically improves positional ordering,
which is exactly what CFID then consolidates.

\subsection{Inference Strategy Ablation}

\begin{table}[H]
\centering
\caption{Inference strategy ablation on the 472-spectrum test set (mean across 3 seeds).
  V1 uses epoch-150 PeakEnc+BYBias weights; V2 adds self-conditioning, cosine schedule,
  AdamW, and 200 epochs.}
\label{tab:ablation}
\begin{tabular}{lccc}
\toprule
\textbf{Inference strategy} & \textbf{Model} & \textbf{AA Recall (\%)} & \textbf{Pep.\ Acc.\ (\%)} \\
\midrule
Argmax (baseline)       & V1 & 74.81 & 36.65 \\
NOVEL\#11: CFID         & V1 & 75.14 & 45.06 \\
NOVEL\#11: CFID + SGIR  & V1 & 75.17 & 44.63 \\
\midrule
Argmax (baseline)       & V2 & 76.21 & 57.84 \\
NOVEL\#9: SGIR only     & V2 & 76.17 & 57.63 \\
NOVEL\#11: CFID         & V2 & 75.93 & 59.60 \\
\rowcolor{gray!15}
NOVEL\#11: CFID + SGIR  & V2 & \textbf{76.00} & \textbf{59.96} \\
NOVEL\#10: spectral rerank & V2 & 73.73 & 53.67 \\
Option B: mass-correct  & V2 & 75.24 & 47.88 \\
\bottomrule
\end{tabular}
\end{table}

In V2, argmax already achieves 76.21\% AA / 57.84\% pep, so CFID adds a smaller
but consistent $+2.11$~pp peptide accuracy gain. CFID+SGIR is the overall best
configuration at 75.99\%~AA / 59.96\%~pep. Spectral reranking and mass correction
both hurt, consistent with V1 findings --- temperature sampling introduces more
noise than the signal recovers at this model quality level.

\paragraph{Ensemble.}
Earlier ablations found that an LSTM+GRU ensemble scored 45.82\% AA recall, well
below the diffusion model alone. Both baselines collapse to repetitive high-frequency
residue predictions on longer sequences, introducing noise rather than complementary
signal when pooled with a stronger model.

\subsection{Training Curves and Gate Confidence}

\begin{figure}[H]
  \centering
  \begin{subfigure}[b]{0.48\textwidth}
    \includegraphics[width=\textwidth]{loss_curves.png}
    \caption{Train/val cross-entropy over 50 epochs (seed 42). Loss decreases
      from 1.94 to 0.99; the train/val gap remains small throughout.}
    \label{fig:loss}
  \end{subfigure}
  \hfill
  \begin{subfigure}[b]{0.48\textwidth}
    \includegraphics[width=\textwidth]{gate_confidence_histogram.png}
    \caption{Gate confidence distribution on the test set (fraction of positions where
      the tight 0.02~Da gate held). Low confidence reflects high model entropy at
      those positions; the 0.1~Da fallback engages so that at least one token remains
      feasible, preventing generation collapse.}
    \label{fig:gate}
  \end{subfigure}
  \caption{Training dynamics and gate behaviour for the diffusion model.}
  \label{fig:training}
\end{figure}

\subsection{ESM-2 Pseudo-Perplexity}

\begin{table}[H]
\centering
\caption{ESM-2 PPL statistics per source. Lower = more consistent with known
  protein sequences.}
\label{tab:esm}
\begin{tabular}{lrrr}
\toprule
\textbf{Source} & \textbf{Mean PPL} & \textbf{Std} & \textbf{n} \\
\midrule
Diffusion   & 31.03 & 10.34 & 472 \\
GRU         & 31.33 & 10.79 & 472 \\
LSTM        & 31.43 & 12.18 & 472 \\
Ground truth & 31.38 & 11.21 & 354 \\
Random      & 35.19 &  7.35 &  50 \\
\bottomrule
\end{tabular}
\end{table}

All three models score near the ground truth mean ($\approx31$), and notably below
random sequences (35.19). The diffusion model has the lowest mean PPL of the three.

\begin{figure}[H]
  \centering
  \begin{subfigure}[b]{0.48\textwidth}
    \includegraphics[width=\textwidth]{figure1_esm2_violin.png}
    \caption{PPL distributions by source. Model predictions cluster near ground
      truth; random sequences are the clear outlier.}
  \end{subfigure}
  \hfill
  \begin{subfigure}[b]{0.48\textwidth}
    \includegraphics[width=\textwidth]{figure2_esm2_heatmap.png}
    \caption{Per-residue PPL heatmap. Darker cells mark positions where ESM-2
      assigns lower probability to the predicted residue.}
  \end{subfigure}
  \caption{ESM-2 pseudo-perplexity across model outputs.}
  \label{fig:esm}
\end{figure}

\subsection{Ensemble Grid Search}

\begin{figure}[H]
  \centering
  \includegraphics[width=0.55\textwidth]{figure3_ensemble_heatmap.png}
  \caption{Validation AA recall across $(\lambda, \gamma)$ pairs. Best result at
    $\lambda=0.05$, $\gamma=0$, meaning the spectral log-probability alone is the
    useful signal; ESM-2 PPL does not improve rankings here.}
  \label{fig:ensemble}
\end{figure}

\subsection{Wastewater Identifications}

\begin{table}[H]
\centering
\caption{De novo PSMs from four wastewater mzML files after 5\% FDR filtering.}
\label{tab:ww}
\begin{tabular}{lccc}
\toprule
\textbf{Sample} & \textbf{Total PSMs} & \textbf{at 5\% FDR} & \textbf{Unique seqs.} \\
\midrule
Sample 1, rep.\ 1 & 4  & 0 & 0 \\
Sample 1, rep.\ 2 & 6  & 0 & 0 \\
Sample 2, rep.\ 1 & 3  & 3 & 3 \\
Sample 2, rep.\ 2 & 3  & 3 & 3 \\
\midrule
\textbf{Total}    & \textbf{16} & \textbf{6} & \textbf{6} \\
\bottomrule
\end{tabular}
\end{table}

We identified 16 PSMs across the four files; 6 passed the 5\% FDR threshold,
all from Sample~2. Identifications from Sample~1 did not survive filtering. The six
high-confidence sequences are listed in Table~\ref{tab:ww_seqs}. All six have ESM-2
PPL below the ground-truth E.~coli mean (31.38), indicating biological plausibility
by the language-model criterion. Since no reference proteome is available for these
communities, BLAST-level validation would require matched metagenome assembly.

\begin{table}[H]
\centering
\caption{Six de novo PSMs passing 5\% FDR in wastewater Sample~2.
  ESM-2 PPL below the ground-truth mean (31.38) indicates biological plausibility.}
\label{tab:ww_seqs}
\begin{tabular}{llcc}
\toprule
\textbf{Sample} & \textbf{Sequence} & \textbf{z} & \textbf{ESM-2 PPL} \\
\midrule
Sample 2, rep.\ 1 & FNDVIPMGEQAINTNEGAYR & 2 & 41.13 \\
Sample 2, rep.\ 1 & NNGNAIGVDLAAIPFVAGDR & 3 & 31.51 \\
Sample 2, rep.\ 1 & GSNYNEVVTLADVTIVQGIR & 3 & 36.63 \\
Sample 2, rep.\ 2 & DLDVEFTALDGASVQVIAYR & 2 & 38.35 \\
Sample 2, rep.\ 2 & ALDNAIDGGQYSFLEVAINR & 2 & 37.87 \\
Sample 2, rep.\ 2 & QLDNNCVYLGATAGVPIAK  & 2 & 37.42 \\
\bottomrule
\end{tabular}
\end{table}

% ─────────────────────────────────────────────────────────────────────────────
\section{Discussion}
% ─────────────────────────────────────────────────────────────────────────────

\paragraph{Why PeakEncoder works.}
The key bottleneck in our earlier model was information loss: a 20,000-bin histogram
discards exact peak positions. The PeakEncoder operates on the raw $(m/z, i)$ pairs,
preserving the sub-Da resolution that distinguishes, e.g., asparagine (114.04~Da)
from isoleucine/leucine (113.08~Da). The $+32$~point improvement over the best
binned-encoder baseline (42.13\% AA) is entirely attributable to this representation.

\paragraph{B/Y-ion bias as physics grounding.}
The learnable scalar $s$ initialised to zero lets the model start as a standard
cross-attention decoder and gradually lean on fragment-ion proximity. This
AlphaFold3-inspired pattern makes the physics injection safe: if the ion signal
were misleading, $s$ would stay near zero; in practice it grows, confirming the
signal is useful.

\paragraph{Self-conditioning closes the feedback loop.}
The single largest gain in V2 is self-conditioning: +14.61~pp peptide accuracy over
V1 CFID (44.99\% → 59.60\%). Without self-conditioning, each denoising step treats
the sequence independently --- the model cannot use committed positions to sharpen
uncertain ones until CFID runs at $t=0$. Self-conditioning brings this feedback into
every reverse step, so by the time CFID runs the sequence is already much more
coherent. The near-zero seed variance ($\pm0.05\%$ AA) indicates the training
procedure is now deterministic enough that results are architecture-driven, not
luck-driven.

\paragraph{Wastewater coverage.}
Six identifications at 5\% FDR is modest. The model was trained exclusively on
clean Orbitrap E.~coli spectra; wastewater acquisitions differ in noise profile,
charge distribution, and peptide composition. Peak-level augmentation (intensity
noise, m/z jitter, random dropout) partially mitigates this mismatch. The recovered
sequences all score below the ground-truth ESM-2 PPL mean, suggesting biological
plausibility; BLAST against NCBI~nr or matched metagenome assemblies would constitute
definitive validation.

\paragraph{Limitations.}
The model is constrained to peptides $\leq\!32$ residues and was trained on Orbitrap
data from a single organism. Generalization to Q-TOF or ion-trap instruments, or to
post-translational modifications beyond carbamidomethyl-C and oxidized-M, has not been
evaluated. The entropy-gate constants (0.02~Da base, scale 2.0) were tuned on E.~coli
validation data and may require re-tuning for other datasets. Finally, we use
mask-predict decoding at inference rather than full iterative reverse diffusion ($T\to0$
with learned re-noising); iterative decoding is expected to improve peptide accuracy
further~\cite{primenovo}.

\paragraph{Reproducibility.}
Code and checkpoints for all three seeds are available at
\url{https://github.com/AkshayRevankarDev/peptide-diffusion}. Training runs on a
single RTX~4090 for approximately 4~hours per seed (200 epochs, batch 256).
All reported numbers are means $\pm$ standard deviations across seeds 0, 1, and~2.

% ─────────────────────────────────────────────────────────────────────────────
\bibliographystyle{unsrt}
\begin{thebibliography}{99}
\small

\bibitem{deepnovo} Tran N, Zhang A, Xin L, et al. De novo peptide sequencing by deep learning. \textit{PNAS} 114(31), 2017.

\bibitem{casanovo} Yilmaz M, Fondrie W, Bittremieux W, et al. De novo mass spectrometry peptide sequencing with a transformer model. \textit{Nat.~Commun.}, 2024.

\bibitem{instanovo} Eloff K, Kalogeropoulos K, Mabona A, et al. InstaNovo enables diffusion-powered de novo peptide sequencing in large-scale proteomics experiments. \textit{Nat.~Mach.~Intell.}~7, 565--579, 2025.

\bibitem{primenovo} Zhang X, et al. $\pi$-PrimeNovo: an accurate and efficient non-autoregressive deep network for de novo peptide sequencing. \textit{Nat.~Commun.}~16, 267, 2025.

\bibitem{ranknovo} Qiu Z, et al. RankNovo: universal biological sequence reranking for de novo peptide sequencing. \textit{ICML}, 2025.

\bibitem{d3pm} Austin J, Johnson D, Ho J, et al. Structured denoising diffusion models in discrete state-spaces. \textit{NeurIPS}, 2021.

\bibitem{esm2} Lin Z, Akin H, Rao R, et al. Evolutionary-scale prediction of atomic-level protein structure with a language model. \textit{Science} 379, 2023.

\bibitem{ofs} Kantroo M, Tauber J, Bhatt D. Pseudo-perplexity in one fell swoop for protein fitness estimation. \textit{PRX Life}, 2024.

\bibitem{novobench} Zhou J, et al. NovoBench: benchmarking de novo sequencing methods in proteomics. \textit{NeurIPS}, 2024.

\bibitem{chen2022} Chen T, Zhang R, Hinton G. Analog Bits: Generating discrete sequences with self-conditioning diffusion models. \textit{ICLR}, 2023.

\bibitem{mdlm} Sahoo SS, et al. Simple and effective masked diffusion language models. \textit{NeurIPS}, 2024.

\end{thebibliography}

\end{document}
